\section{Методы испытаний}

\subsection{Проверка запуска и справочной информации}

Испытание проводится путем запуска программы с параметром вывода справочной
информации:
\begin{verbatim}
schengen_main --help
\end{verbatim}

Положительным результатом испытания считается завершение программы без ошибки и
вывод сведений о поддерживаемых параметрах запуска.

\subsection{Проверка вывода служебной информации}

Испытание проводится путем последовательного запуска команд:
\begin{verbatim}
schengen_main --list-tables
schengen_main --list-formats
\end{verbatim}

Положительным результатом испытания считается вывод перечня доступных таблиц
TPC-H и перечня форматов вывода, предусмотренных текущей сборкой программы.

\subsection{Проверка соответствия эталонному генератору TPC-H}

Испытание проводится путем сопоставления результатов генерации с контрольными
таблицами, полученными исходным генератором набора TPC-H. Для указанной
проверки используются контрольные примеры для таблиц \texttt{region},
\texttt{nation}, \texttt{supplier}, \texttt{part}, \texttt{partsupp},
\texttt{customer}, \texttt{orders} и \texttt{lineitem}.

Проверка может выполняться посредством автоматизированных тестов генераторов
таблиц либо иным способом построчного сопоставления контрольных значений с
эталонными файлами \texttt{*.tbl}, полученными оригинальным TPC-H генератором.

Положительным результатом испытания считается совпадение с эталонными данными
по всем контролируемым таблицам и отсутствие расхождений в составе и значениях
сформированных строк.

\subsection{Проверка генерации полного набора таблиц}

Испытание проводится путем запуска программы в режиме генерации полного набора
таблиц:
\begin{verbatim}
schengen_main --scale-factor 1 --output-format parquet --output-path ./out
\end{verbatim}

В процессе испытания проверяется факт создания выходного каталога и наличие
результатов генерации, соответствующих таблицам набора TPC-H.

Положительным результатом испытания считается завершение программы с кодом
\texttt{0} и формирование в каталоге вывода данных полного набора таблиц.

\subsection{Проверка генерации выбранных таблиц}

Испытание проводится путем запуска программы в режиме генерации выбранного
подмножества таблиц:
\begin{verbatim}
schengen_main \
  --scale-factor 1 \
  --output-format parquet \
  --output-path ./out_selected \
  --table customer orders lineitem
\end{verbatim}

Положительным результатом испытания считается формирование только указанных
оператором таблиц при отсутствии выходных данных для остальных таблиц набора.

\subsection{Проверка параметров записи}

Испытание проводится путем запуска программы с параметрами батчирования и
записи, предусмотренными текущей сборкой. В качестве проверочного примера может
использоваться следующая команда:
\begin{verbatim}
schengen_main \
  --scale-factor 1 \
  --output-format parquet \
  --output-path ./out_tuned \
  --batch-rows 10000 \
  --write-strategy auto \
  --write-workers 0
\end{verbatim}

Положительным результатом испытания считается прием указанных параметров,
выполнение генерации без сбоев и сохранение выходных данных в заданном каталоге
вывода.

\subsection{Проверка устойчивой работы при больших значениях масштаба генерации}

Испытание проводится путем запуска программы при значении масштаба генерации,
превышающем 100. В качестве проверочного примера может использоваться
следующая команда:
\begin{verbatim}
schengen_main \
  --scale-factor 300 \
  --output-format parquet \
  --output-path ./out_sf300 \
  --table orders lineitem
\end{verbatim}

В процессе испытания контролируется отсутствие аварийного завершения,
диагностики, связанной с переполнением внутренних диапазонов генерации,
некорректного формирования ключевых значений и иных ошибок, обусловленных
увеличением масштаба генерируемого набора данных.

Положительным результатом испытания считается завершение программы с корректным
формированием выходных данных для выбранных таблиц и отсутствие признаков
неустойчивой работы при большом значении \texttt{scale factor}.

\subsection{Проверка диагностирования ошибочных параметров}

Испытание проводится путем запуска программы с заведомо некорректными
параметрами. В качестве проверочных примеров могут использоваться команды:
\begin{verbatim}
schengen_main --scale-factor 0 --output-format parquet --output-path ./out
schengen_main --scale-factor 1 --output-format unknown --output-path ./out
\end{verbatim}

Положительным результатом испытания считается завершение программы с ненулевым
кодом возврата и выдача диагностических сообщений, указывающих на причину
отказа.

\subsection{Проверка программной документации}

Испытание проводится путем сопоставления документации, предъявленной на
испытания, с фактическим поведением программы при выполнении проверок,
предусмотренных настоящим документом.

Положительным результатом испытания считается полнота состава документации и
отсутствие противоречий между сведениями о назначении программы, параметрах
запуска, результатах работы и фактическими результатами проведенных испытаний.
